\section{Программирование}

\subsection{Реализация практических работ}

В первой практической работе были реализованы функции для работы со списками и пользовательскими типами. Вторая работа содержит проекты \mintinline{text}{MyEvolProject} и \mintinline{text}{myParser}. Третья работа содержит проекты обработки изображения и поиска в лабиринте. Четвёртая работа реализует интерактивную игру по лабиринту. Пятая работа содержит функции и QuickCheck-тесты.

Функция запуска игры в четвёртой работе последовательно раскрывает стек трансформеров и запускает итоговое \mintinline{haskell}{IO}-действие:

\captionof{listing}{Запуск игры и терминирование трансформеров}
\label{lst:run-game}
\begin{minted}{haskell}
runGame roomPairs game =
  runMyStateT (runWriterT (runReaderT game (buildMaze roomPairs))) "start"
\end{minted}

Сначала \mintinline{haskell}{runReaderT} получает окружение --- построенный лабиринт. Затем \mintinline{haskell}{runWriterT} раскрывает слой лога. После этого \mintinline{haskell}{runMyStateT} получает начальное состояние \mintinline{haskell}{"start"} и возвращает действие в \mintinline{haskell}{IO}.

В пятой работе генераторы и свойства QuickCheck позволяют проверять функции на большом количестве автоматически сгенерированных входных данных.

\captionof{listing}{Генераторы для QuickCheck}
\label{lst:quickcheck-generators}
\begin{minted}{haskell}
genWord :: Gen String
genWord = listOf1 $ elements ['a' .. 'z']

genSentence :: Gen String
genSentence =
  choose (0, 8) >>= \n ->
    replicateM n genWord >>= \wordsList ->
      return $ intercalate " " wordsList

genMod :: Gen Int
genMod = choose (1, 1000)
\end{minted}

\subsection{Использование нейросетей}

Использование нейросетей фиксировалось в файлах практических работ. В первой работе модель указана в комментариях к исходному файлу. Во второй работе сведения о модели вынесены в README и комментарии к файлам практики. В третьей работе добавлен файл с историей запросов к LLM для проекта \mintinline{text}{glitcher}. В четвёртой работе был создан файл переписки с LLM в проекте \mintinline{text}{mazeGame}; после завершения работы лог больше не дополнялся. В пятой работе история взаимодействия сохранена в \mintinline{text}{practice_05/qc/llm.txt}.

\subsection{История коммитов}

В таблице~\ref{table:commits} приведены основные коммиты по практическим работам.

{\small
\begin{longtblr}[
    caption = {Основные коммиты по практическим работам},
    label = {table:commits},
]{
    colspec = {|X[0.8,c]|X[1.1,c]|X[5,l]|},
    row{1} = {font=\bfseries},
    hlines,
    vlines,
}
Практика & Коммит & Сообщение \\
1 & \mintinline{text}{f0752bd} & Добавлены функции myZipSave \\
1 & \mintinline{text}{5baf79c} & Добавлены функции myUnzipSave \\
1 & \mintinline{text}{78c0460} & Добавлены функции myReverse, myFoldl1, myFoldr1 \\
1 & \mintinline{text}{c575c1c} & Добавлена функция myMaybe \\
1 & \mintinline{text}{5b5342c} & Ответ на вопрос на защите первой лабораторной работы \\
2 & \mintinline{text}{49f3596} & Добавлен файл MyEvolModule; добавлена сборка через cabal \\
2 & \mintinline{text}{bd2c465} & Добавлены представители классов типов Show, Eq, Enum для MyEvolution \\
2 & \mintinline{text}{75ef59e} & Добавлен MyEvolution' \\
2 & \mintinline{text}{17cca30} & Инициализирован stack-проект MyParser \\
2 & \mintinline{text}{dc5e0eb} & Добавлены парсеры, реализованные с использованием библиотек \\
3 & \mintinline{text}{d04a59e} & Инициализирован stack-проект glitcher \\
3 & \mintinline{text}{675192e} & Добавлен вызов нескольких глитчей по порядку в Main \\
3 & \mintinline{text}{aac32e4} & Инициализирован stack-проект с лабиринтом \\
3 & \mintinline{text}{b766d6e} & Добавлены функции поиска в лабиринте \\
3 & \mintinline{text}{88cc5cd} & Поправлен вывод в файлы в 3 практике \\
4 & \mintinline{text}{e4678a9} & Инициализирован stack-проект mazeGame для 4 работы \\
4 & \mintinline{text}{94e6e47} & Добавлен лабиринт главного здания для 4 практики \\
4 & \mintinline{text}{1f73730} & Добавлена функция чтения лабиринта из файла \\
4 & \mintinline{text}{bba7573} & Добавлен трансформер MyStateT \\
4 & \mintinline{text}{95f97e2} & Добавлены функции runGame, gameLoop и вспомогательные функции \\
4 & \mintinline{text}{99277ac} & Revert: убран переход назад из 4 практики \\
4 & \mintinline{text}{515420e} & Изменён порядок трансформеров в типе Game \\
5 & \mintinline{text}{7bf6419} & Инициализирован stack-проект qc в практике №5 \\
5 & \mintinline{text}{ad439ae} & Добавлены функции addMod и reverseWords \\
5 & \mintinline{text}{928bdbe} & Добавлены QuickCheck-тесты для функций 5 практики \\
5 & \mintinline{text}{462c309} & Обновлены QuickCheck-тесты и добавлена история llm \\
5 & \mintinline{text}{c7d26e5} & Тесты переписаны через forAll \\
\end{longtblr}
}

\newpage
